\documentclass[10pt,a4paper]{amsart}
\usepackage[margin=19mm]{geometry}
\usepackage{amsmath,amssymb,mathtools}
\usepackage[scr=rsfs,cal=euler]{mathalfa}
\usepackage{xfrac}
\usepackage[unicode,hidelinks,bookmarks=false]{hyperref}
\newcommand{\ie}{i.e.\ }
\newcommand{\cf}{cf.\ }
\newcommand{\resp}{respectively\ }
\newcommand{\Subsection}{\subsection}
\providecommand{\nobreakdash}{\nobreak\hbox{-}}
\setlength{\emergencystretch}{2em}
\multlinegap0pt
\usepackage{bm}
\usepackage{bbm}
\usepackage{dsfont}
\usepackage{xspace}
\usepackage{cancel}
\usepackage{mathtools}
\usepackage[shortlabels]{enumitem}
\usepackage{amsthm}
\makeatletter
\@ifundefined{theorem}{\newtheorem{theorem}{Theorem}}{}
\@ifundefined{lemma}{\newtheorem{lemma}[theorem]{Lemma}}{}
\makeatother
\makeatletter
\expandafter\ifx\csname case\endcsname\relax
\newenvironment{case}[2][Case]{%
  \trivlist \item[\hskip\labelsep{\itshape #1 #2.}]\begin{em}}{
  \end{em}\endtrivlist}
\fi
\expandafter\ifx\csname subcase\endcsname\relax
\newenvironment{subcase}[2][Subcase]{%
  \trivlist \item[\hskip\labelsep{\itshape #1 #2.}]\begin{em}}{
  \end{em}\endtrivlist}
\fi
\makeatother
\title[Characterization of the tree cycles with minimum positive en]{Characterization of the tree cycles with minimum positive entropy for any period}
\author{David Juher, Francesc Ma\~nosas, David Rojas}
\date{Source: arXiv:2310.14862v1 (CC BY 4.0)}
\begin{document}
\maketitle

\noindent\emph{Source and licence.} Characterization of the tree cycles with minimum positive entropy for any period. Version arXiv:2310.14862v1 (CC BY 4.0), distributed under the Creative Commons Attribution 4.0 International licence, as recorded in the arXiv licence metadata for this exact version and shown on the abstract page https://arxiv.org/abs/2310.14862v1. Full rights notice: licenses/arxiv-2310.14862-CC-BY-4.0.txt.\par\smallskip

\noindent\textit{Editorial scope.} Counted unit: the Lemma whose source label is \texttt{isi2} in the article (source0.tex lines 799--801, source label \texttt{isi2}), delivered together with the article's own complete original proof of it (source0.tex lines 802--820, one \texttt{proof} environment of 19 lines). No mathematics is written, repaired or re-derived: statement and proof are the article's own text line for line. Required context retained uncounted: none. Every definition, display and label of those lines is retained unchanged. Omitted from this package and not redistributed: 1--798, 821--1798. Nothing inside a retained excerpt is omitted or altered: every retained body is delivered exactly as the article's lines are written, so no omission receipt of any kind accompanies this package. Citations: the retained excerpts carry no citation command at all, so no bibliography entry of the article is transcribed and no reference is redistributed. Reference adaptation: none was needed and none was made; every reference inside the delivered unit resolves inside it, so no unresolved reference or citation is produced. Printed slips kept literal: no printed word, symbol, hypothesis or number of the counted statement or of its proof was altered. Layout and class: the article's own class is \texttt{amsart} with packages including inputenc, graphicx, color, easybmat, amssymb, amsmath, enumerate; the wrapper is the lane's pinned \texttt{amsart} editorial wrapper at 10pt on A4, which restates the theorem-like environments the retained text needs and adds the article's own macro definitions that the retained text needs (none), with extra wrapper packages (bm, bbm, dsfont, xspace, cancel, mathtools, enumitem). The delivered numbering is the unit's own. Assets: no retained span contains a figure environment, a graphics-inclusion command, a PDF-page inclusion, a table or a TikZ picture; no image file is copied into this package, the package declares no asset dependency and no image file is redistributed with it. Licence: version arXiv:2310.14862v1 of this article is distributed under the Creative Commons Attribution 4.0 International licence, as recorded in the arXiv licence metadata for this exact version and shown on the abstract page https://arxiv.org/abs/2310.14862v1. A full rights notice is delivered at licenses/arxiv-2310.14862-CC-BY-4.0.txt. Characters: every retained excerpt is pure ASCII, so the delivered engine, XeLaTeX with its default Latin Modern text font, reports no missing character.
\begin{lemma}\label{isi2}
Let $\mathcal{P}$ be a zero entropy periodic pattern with a maximal separated structure $P_0\cup P_1\cup\ldots\cup P_{p-1}$ of trivial blocks. Let $\mathcal{C}$ be the combinatorial collapse of $\mathcal{P}$. If $0\le i,j,k<p$ are three points of $\mathcal{C}$ such that $\{j,k\}$ is contained in a single $i$-branch of $\mathcal{C}$, then $P_j\cup P_k$ is contained in a single $i$-branch of $\mathcal{P}$.
\end{lemma}

\begin{proof}
Assume by way of contradiction that $P_j$ and $P_k$ are respectively contained in two different $i$-branches of $\mathcal{P}$. Then, for some $N\ge 2$ there exist $N+1$ different points of $\mathcal{P}$, $x_0,x_1,\ldots,x_N$, such that:
\begin{enumerate}
\item $x_0=j$ and $x_N=k$
\item $x_n$ is inner for all $0<n<N$
\item $\{x_n,x_{n+1}\}$ is contained in a discrete component of $\mathcal{P}$ for $0\le n<N$
\item $x_m=i$ for some $1\le m<N$
\end{enumerate}
Intuitively, the above ordered sequence of points accounts for all points of $\mathcal{P}$ successively met in the shortest path going from $j$ to $k$. The assumption that $i$ separates $j$ from $k$ is imposed by property (d).

Consider now, for any point $x_n$ of the above sequence, the collapse of the trivial block of $\mathcal{P}$ containing $x_n$. It is a point of $\mathcal{C}$ that we denote by $y_n$. Note that $y_0=j$, $y_m=i$ and $y_N=k$. Observe also that, for a pair of consecutive points $x_n,x_{n+1}$, it may happen that $\{x_n,x_{n+1}\}$ is contained in a block. In this case, since the blocks are trivial, $\{x_n,x_{n+1},x_{n+2}\}$ is not contained in a block. Therefore, $y_n=y_{n+1}\ne y_{n+2}$. On the other hand, if $\{x_n,x_{n+1}\}$ is not contained in a block, by the definition of the combinatorial collapse, $\{y_n,y_{n+1}\}$ is a binary set contained in a discrete component of $\mathcal{C}$. This observations lead to the existence of a sequence $z_0,z_1,\ldots,z_{M}$ of $M+1\le N+1$ points of $\mathcal{C}$ such that
\begin{enumerate}[(a')]
\item $z_0=j$ and $z_M=k$
\item $z_n$ is inner for all $0<n<M$
\item $\{z_n,z_{n+1}\}$ is contained in a discrete component of $\mathcal{C}$ for $0\le n<M$
\item $x_{m'}=i$ for some $1\le m'<M$
\end{enumerate}
By property (d'), $j$ and $k$ belong to different $i$-branches in $\mathcal{C}$, in contradiction with the hypothesis of the lemma.
\end{proof}

\end{document}
